\chapter{Unique Individuals}\label{ch:g9:unique-individuals}

Seven billion people, and no two alike --- not in a family,
not in a schoolyard, not in all of history, identical twins
half-excepted. Three chapters have now assembled the whole
machinery of that fact. This short capstone chapter runs the
argument once, end to end, and banks its consequences: for
the \omterm{def:g5:biodiversity:species}{species}' variety, for how we read ``genetic''
headlines, and for what uniqueness does and does not mean.

\section{The uniqueness theorem}

\begin{proposition}[Why no two people match]\label{prop:g9:unique-individuals:theorem}
Assemble the three mechanisms:
\begin{enumerate}
  \item \emph{the shuffle of the halving}: each parent's
        \omterm{def:g8:sexual-reproduction:gametes}{gametes} carry one partner per pair, dealt
        independently --- $2^{23}$ half-decks per parent
        (\cref{rem:g9:chromosomes:whyhalve}), and in truth
        far more, since paired partners also \emph{swap
        segments} during the halving division, cutting the
        deck below the level of whole \omterm{def:g9:chromosomes:chromosomes}{chromosomes};
  \item \emph{the lottery of \omterm{def:g5:flower-to-fruit:steps}{fertilization}}: which sperm
        of the millions (\cref{ex:g8:reproductive-systems:sperm})
        meets which month's \omterm{def:g2:animals-grow-up:born}{egg} multiplies the two spaces
        together;
  \item \emph{\omterm{prop:g9:genes-and-dna:explains}{mutation}'s trickle}: every generation adds a
        few fresh spellings of its own
        (\cref{prop:g9:genes-and-dna:explains}).
\end{enumerate}
The combination space dwarfs the number of humans who have
ever lived: every \omterm{def:g5:flower-to-fruit:steps}{fertilization} deals a deck that has
never been dealt before, and never will be again.
\end{proposition}
\admitted

\begin{example}[The twins, and the exception's limits]\label{ex:g9:unique-individuals:twins}
Identical twins share one deal
(\cref{exo:g8:from-egg-to-birth:10}) --- and even they
diverge: separate lives write different acquired \omterm{def:g9:heredity-and-traits:traits}{traits}
(\cref{prop:g9:heredity-and-traits:sorting}), different
\omterm{def:g8:nervous-communication:synapse}{synapse} histories
(\cref{ex:g8:nervous-communication:juggler}), even a few
different late \omterm{prop:g9:genes-and-dna:explains}{mutations} in each body's own \omterm{def:g5:first-look-at-cells:cell}{cells}. The one
crack in uniqueness seals itself: after the deal, living
differentiates.
\end{example}

\begin{remark}[Unique, and mostly alike]\label{rem:g9:unique-individuals:alike}
Hold both truths: any two humans' \omterm{def:g9:genes-and-dna:dna}{DNA} texts agree over
$99$ percent --- the shared build of
\cref{prop:g5:first-look-at-cells:principle}'s one \omterm{def:g5:biodiversity:species}{species}
--- and the disagreements, scattered through three billion
letters, still suffice for strict uniqueness. ``All
different'' and ``all kin'' are the same fact at two
magnifications
(\cref{met:g9:genes-and-dna:levels}); every use ever made
of the first truth against the second was bad \omterm{def:g1:living-or-not:living}{biology}
before it was bad ethics.
\end{remark}

\section{Uniqueness at work}

\begin{example}[DNA identification]\label{ex:g9:unique-individuals:forensics}
Uniqueness is checkable, and courts check it: enough of a
person's variable positions, read from a trace of \omterm{def:g5:first-look-at-cells:cell}{cells},
match one individual (identical twins aside) out of
humanity. The same comparison, run at family
magnification, resolves parentage --- the child's every
variable position must be dealt from its two parents'
decks. The album's logic
(\cref{ch:g9:heredity-and-traits}) has become an
instrument.
\end{example}

\begin{example}[Transfusion and transplantation]\label{ex:g9:unique-individuals:medicine}
Medicine negotiates uniqueness daily. \omterm{prop:g4:journey-of-food:absorb}{Blood} groups
(\cref{ex:g9:genes-and-dna:bloodgroups}) are the easy
case --- four broad classes, matchable from the register.
Organ transplants meet the hard case: the immune system
(\cref{ch:g9:immune-defenses} explains the machinery)
reads finer-grained self-markers, unique enough that
donors are searched for by family first --- the closer the
\omterm{prop:g6:classification-kinship:kinship}{kinship}, the nearer the deal.
\end{example}

\begin{proposition}[Variety is the species' reserve, banked]\label{prop:g9:unique-individuals:reserve}
\Cref{rem:g8:sexual-reproduction:reserve} promised the
shuffle's point; the mechanisms now cash it. A \omterm{def:g5:biodiversity:species}{species} of
unique individuals holds, scattered among them, \omterm{def:g9:genes-and-dna:gene}{alleles}
for every occasion --- the resistance some carry to this
disease, the tolerance others carry to that stress
(\cref{ex:g8:sexual-reproduction:bets}'s orchard, at
\omterm{def:g5:biodiversity:species}{species} scale). When \omterm{def:g6:exploring-our-environment:conditions}{conditions} turn, the checks of
\cref{prop:g8:reproduction-and-environment:limits} fall
unevenly across that variety --- and what that unevenness
does, generation after generation, is precisely the
subject the next two chapters open: the changing of
\omterm{def:g5:biodiversity:species}{species}.
\end{proposition}
\admitted

\begin{method}[Auditing a ``genetic'' claim]\label{met:g9:unique-individuals:claims}
For any headline about \omterm{def:g9:genes-and-dna:gene}{genes} and people:
\begin{enumerate}
  \item level check (\cref{met:g9:genes-and-dna:levels}):
        is the claim about a \omterm{def:g9:genes-and-dna:gene}{gene}, a deck, or a family
        pattern?
  \item weave check
        (\cref{prop:g9:heredity-and-traits:sorting}): does
        it respect the inherited-range-lived-in weave, or
        pretend one \omterm{def:g9:genes-and-dna:gene}{gene} equals one destiny?
  \item variety check: does it treat a \omterm{def:g8:reproduction-and-environment:population}{population}'s spread
        as the reserve it is, or as deviations from a
        ``normal'' deck that does not exist?
  \item twin check: would it survive the identical twins
        who share the deal and differ anyway?
\end{enumerate}
\end{method}

\section{Exercises}

\begin{exercise}[$\star$]\label{exo:g9:unique-individuals:1}
Name the three mechanisms of the uniqueness theorem.
\end{exercise}

\begin{exercise}[$\star$]\label{exo:g9:unique-individuals:2}
What does segment-swapping during the halving add to the
shuffle?
\end{exercise}

\begin{exercise}[$\star$]\label{exo:g9:unique-individuals:3}
State both halves of \cref{rem:g9:unique-individuals:alike}
--- the percentage and the uniqueness --- in one
sentence.
\end{exercise}

\begin{exercise}[$\star$]\label{exo:g9:unique-individuals:4}
How do identical twins end up distinguishable after all?
\end{exercise}

\begin{exercise}[$\star$]\label{exo:g9:unique-individuals:5}
What two questions does \omterm{def:g9:genes-and-dna:dna}{DNA} comparison answer for a
court?
\end{exercise}

\begin{exercise}[$\star$]\label{exo:g9:unique-individuals:6}
Why are transplant donors searched for by family first?
\end{exercise}

\begin{exercise}[$\star\star$]\label{exo:g9:unique-individuals:7}
Multiply the argument: combine the two parents'
half-deck spaces and the \omterm{def:g5:flower-to-fruit:steps}{fertilization} lottery into the
``never dealt before'' conclusion.
\end{exercise}

\begin{exercise}[$\star\star$]\label{exo:g9:unique-individuals:8}
Cash \cref{prop:g9:unique-individuals:reserve} on the
vineyard of \cref{exo:g8:sexual-reproduction:10}: what
did the clones lack that a seed-grown \omterm{def:g8:reproduction-and-environment:population}{population} holds?
\end{exercise}

\begin{exercise}[$\star\star$]\label{exo:g9:unique-individuals:9}
Run \cref{met:g9:unique-individuals:claims} on: ``a \omterm{def:g9:genes-and-dna:gene}{gene}
for school success discovered''.
\end{exercise}

\begin{exercise}[$\star\star$]\label{exo:g9:unique-individuals:10}
Why is ``the normal human deck'' a phrase without a
referent? Answer with the reserve and the theorem.
\end{exercise}

\begin{exercise}[$\star\star$]\label{exo:g9:unique-individuals:11}
A parentage test clears a doubted father wrongly accused
by \cref{exo:g9:heredity-and-traits:11}'s neighbor.
Explain what the test checks, position by position.
\end{exercise}

\begin{exercise}[$\star\star\star$]\label{exo:g9:unique-individuals:12}
``All different, all kin, and both by the same
mechanism.'' Write the capstone paragraph: shuffle,
\omterm{prop:g9:genes-and-dna:explains}{mutation} and shared code, each placed.
\end{exercise}

\section{Problem: The Cold-Case Seminar}

\begin{problem}[{Weekend problem --- one hair, three
questions, thirty years}]\label{pb:g9:unique-individuals:1}
A (fictional) cold-case seminar: a thirty-year-old
unsolved burglary, one preserved hair with its \omterm{def:g1:plants-around-us:parts}{root}
\omterm{def:g5:first-look-at-cells:cell}{cells}, and modern comparison methods. The seminar walks
students through what \omterm{def:g9:genes-and-dna:dna}{DNA} identification can and cannot
say.

\medskip\noindent\textbf{Part I --- What the hair
holds.}
\begin{enumerate}
  \item Why must the hair's \emph{\omterm{def:g1:plants-around-us:parts}{root}} be preserved for
        a full reading (\cref{def:g6:cell-unit-of-life:parts}
        knows where \omterm{def:g9:genes-and-dna:dna}{DNA} lives)?
  \item The lab reads a set of highly variable
        positions. Why variable ones --- what would the
        99-percent-shared positions establish?
  \item The profile matches a suspect at every position.
        State what the match asserts, and its one
        systematic exception
        (\cref{ex:g9:unique-individuals:twins}).
  \item The defense notes the suspect has an identical
        twin abroad. What must the court now do, and
        why?
\end{enumerate}

\medskip\noindent\textbf{Part II --- The family
questions.}
\begin{enumerate}[resume]
  \item An earlier, partial profile had pointed to ``a
        close relative of the K family''. Explain how a
        \emph{partial} match reads as \omterm{prop:g6:classification-kinship:kinship}{kinship}
        (\cref{prop:g9:unique-individuals:theorem}'s
        dealing, run backward).
  \item A juror asks whether the profile reveals the
        suspect's \omterm{def:g9:heredity-and-traits:traits}{traits} --- looks, temperament,
        ``criminal \omterm{def:g9:genes-and-dna:gene}{genes}''. Run
        \cref{met:g9:unique-individuals:claims} on the
        question, check by check.
  \item Another juror asks whether the match could be
        ``once in this city'' rather than once in
        humanity. What sets the number of positions
        read?
  \item The seminar leader warns: ``the \omterm{def:g1:living-or-not:living}{biology} says
        whose \omterm{def:g5:first-look-at-cells:cell}{cells}; the case must still say how they
        got there.'' Distinguish the two claims.
\end{enumerate}

\medskip\noindent\textbf{Part III --- The seminar's
close.}
\begin{enumerate}[resume]
  \item Summarize the instrument in one sentence: what
        two chapters' mechanisms
        (\cref{ch:g9:chromosomes},
        \cref{ch:g9:genes-and-dna}) make a profile
        unique and a family readable.
  \item Why does the same instrument work on a
        thirty-year-old hair? (What property of \omterm{def:g9:genes-and-dna:dna}{DNA} ---
        \cref{prop:g9:genes-and-dna:explains}'s copying
        aside --- is being relied on: stability.)
  \item The seminar ends on ethics: list two uses of
        uniqueness this chapter shows serving people,
        and the historical misuse
        \cref{rem:g9:unique-individuals:alike} warns
        against.
  \item Close with the theorem restated for the
        courtroom: what has never happened twice, and
        why.
\end{enumerate}
\end{problem}
